\section{考察}
\subsection{予備実験から分かること}
予備実験では，同じボックスを用いて物体認識モデル自身のクラスと文章・画像埋め込みモデルによるクラスを比較したため，両者の差は全てクラスの推定に起因する．
COCOの80クラスの範囲では，COCOで学習した物体認識モデルのラベルの一致率は0.95以上であり，ゼロショット分類の0.56〜0.58を大きく上回った．
したがって，学習時のクラスの範囲では物体認識モデル自身のクラスを用いるべきであり，文章・画像埋め込みモデルの利点は，学習時に無い物体と語彙を追加の学習なしに扱える点にある．
また，WordNetの上位語から1段ずつ下位語を選ぶ方法では，役割を表す下位語に進んで細分化が止まる誤りと，最初の段階の誤りが以降の全ての段階に伝播する誤りが見られた．
提案手法は，前者を人の役割の語の除外と語義を区別する文章で，後者を語彙全体の確率の上位語への集約で緩和するものである．

\subsection{上位語への集約の効果と限界}
4.5節から，上位語への集約による命名は，同じ領域に対する命名の正しさを25.9\%から43.0\%に改善した．
これは，表\ref{tab_diag}のように学習時に無い物体に対する1位の単語の正解率が低くても，正解の語とその近くの語の確率が共通の上位語に集まるためである．
一方，改善の多くは粗い語による正解であり，LVISのカテゴリまで正しい出力のみを数えるAP$ _{\mathrm{exact}} $ではflatが上回った．
すなわち，提案手法は細かい語の正解を増やす方法ではなく，誤った細かい語を正しい粗い語に置き換える方法であり，その効果は階層的な指標でのみ現れる．
また，評価の語彙を事前に与えた場合には集約の効果は無かった．
語彙が小さく正解の語を必ず含む場合には，1位の単語の正解率が高く，粗い語に止まる利点が小さいためと考えられる．
したがって，提案手法が有効なのは，評価の語彙を事前に知らず，学習時に無い物体を大きな語彙から命名する場合である．
細かい語の正解を増やすには，ST-LやFTとの比較から分かるように，領域の特徴の分類能力そのものを高める必要がある．

\subsection{実用上の水準と誤りの要因}
提案手法の出力は，そのままでは人手のアノテーションの代わりにはならない．
STを用いて検証用の画像で選んだ閾値を用いると，閾値以上の加えた領域の出力のうち正しく命名されたものは16.6\%（表\ref{tab_errors}）であり，COCO外の物体のうち命名できたものは3.2\%であった．
一方，(\ref{eq_score})のスコアはテスト用の画像で正解率とほぼ一致したため，スコアが高い出力のみを候補として人に示し，確認させる用途には利用できる．
4.7節の誤りの分解から，加えた領域の出力の誤りの大部分は命名ではなく物体の一部や背景に対するものであり，GTと重なる領域のみを残すとAP$ _{\mathrm{strict}} $は0.197から0.265に向上した．
したがって，次に改善すべきは命名ではなく，領域が物体の全体であるかの判定である．
なお，LVISのアノテーションは網羅的でないため，GTと対応しない出力の全てが誤りとは限らず，$ hF_{\mathrm{fed}} $と$ hF_{\mathrm{strict}} $，およびAP$ _{\mathrm{fed}} $とAP$ _{\mathrm{strict}} $は真の値の上限と下限に近いと考えられる．
また，スコアの重みはLVISの検証用の画像で求めたため，LVISと異なる分布の画像では校正が崩れる可能性がある．

\subsection{今後の課題}
第一に，領域が物体の全体であるかを判定し，物体の一部や背景の領域を棄却することが挙げられる．
\add{4.7節のように，ボックスの包含関係や大きさ，WordNetの部分と全体の関係のみでは効果が小さかったため，領域の周囲を含めて切り抜いた画像の特徴の変化等，画像の内容に基づく判定が必要である．}
第二に，学習時に無い物体に対する領域の特徴の分類能力の向上が挙げられ，より規模の大きい文章・画像埋め込みモデルの特徴と組み合わせる方法が考えられる．
第三に，LVISのアノテーションの網羅性の不足を補うため，GTと対応しない出力を人手で確認し，真の適合率を求めることが挙げられる．
